\section{Amenazas a la validez}
\label{sec:amenazas}

\subsection{Identificación y selección de estudios}
\label{sec:am-seleccion}
% PRISMA ítem 23c (limitaciones del proceso de revisión)

La búsqueda abarcó \nBases{} fuentes, con ACM sin filtro de fecha y arXiv consultado por API con filtro local, por lo que la cobertura puede ser incompleta. No se consultaron DBLP ni Web of Science, y las cadenas estaban en inglés. La sensibilidad fue menor para los estudios de un turno (\nRefUTHallados{} de \nRefUT{}; IC 95~\% de \sensUTInf{} a \sensUTSup~\%) que para los multiturno (\nRefMTHallados{} de \nRefMT{}). Si se omitieron más estudios de un turno, el denominador de D4 se subestima y la proporción multiturno se sobreestima, de modo que el sesgo iría en contra de la lectura de minoría. Además, la IA determinó la elegibilidad del conjunto de referencia, lo que puede sobreestimar la sensibilidad por circularidad; por eso el límite inferior multiturno de \refMTICinf~\% es una cota optimista, y no se dispuso de un conjunto de referencia independiente. La contribución central se juzgó por título y resumen, con una clasificación en parte subjetiva; se aplicó la misma regla en ambos escenarios y la duda avanzó a texto completo.

El cribado lo hicieron revisores de IA de la misma familia, con una única revisora humana. Sus errores pueden estar correlacionados, y su acuerdo no equivale a una revisión independiente. La relectura ciega de las exclusiones conjuntas avanzó \nQCAvanzados{} registros, de los que se incluyeron \nQCIncluidos{}: los revisores A y B juntos omitieron estudios elegibles. De las \nConjuntas{} exclusiones conjuntas, \nNoReleidas{} no se releyeron; en el estrato aleatorio, \nRescAleat{} de \nAleatRel{} se incluyeron finalmente (\tasaRescate~\%, IC 95~\% de \tasaRescateInf{} a \tasaRescateSup~\%), lo que, extrapolado, equivale a \omisEsp{} estudios omitidos esperados (IC 95~\% de \omisEspInf{} a \omisEspSup{}). Como la relectura cubrió solo una muestra (\nQC{} registros), pueden persistir omisiones no detectadas. La verificación humana siguió un diseño reducido, con \nPiloto{} registros y sin muestra reservada, por lo que las estimaciones de sensibilidad son imprecisas. La crítica del protocolo la hizo un agente de la misma familia y no sustituye una revisión por pares independiente.
% fuente: revision_4f.json (A5_omisiones, A8_sensibilidad)

En texto completo se usaron extractos y modelos distintos. Codex coincidió con el texto completo en \nCalibCodex{} de \nCalibN{} informes de calibración, frente a \nCalibClaude{} de \nCalibN{} de Claude, por lo que la parte del Nivel 1 que revisó puede diferir en sensibilidad. La auditoría confirmó \nAuditOK{} de \nAuditN{} inclusiones de Codex, pero no examinó sus exclusiones.

IEEE Xplore aportó pocos estudios incluidos en proporción a sus registros (\cref{fig:busqueda_fuentes}), y \nNoRecuperados{} informes del brazo de bases no se recuperaron, sin contactar a los autores. Si la disponibilidad depende de la editorial, los estudios de acceso restringido quedan infrarrepresentados. La señal de los informes no recuperados se asignó por título y resumen, de modo que es solo una aproximación al escenario. Las sensibilidades (una cota extrema y una estimación probabilística) acotan el efecto sobre D4, pero no lo eliminan.

En cuanto a la validez de constructo, las preguntas PI1–PI4 y la taxonomía derivan de la propuesta de tesis de la autora, por lo que pueden privilegiar los componentes que esa propuesta considera relevantes; la crítica del protocolo ya señaló este riesgo de sesgo de confirmación. Se aplicaron varias mitigaciones. El brazo multiturno no se restringió a esos componentes: era elegible cualquier estudio que abordara al menos una de PI1--PI5. Las clases inductivas «otros componentes» y las categorías \texttt{OTRA} se consolidaron y se reportan. El criterio de reformulación (M8) se fijó antes de examinar los datos como regla que obligaría a replantear la brecha de partida, y se reporta con independencia del resultado. El conflicto de interés se declara de forma explícita. Persiste un riesgo residual: las categorías de análisis y la agenda de investigación pueden seguir enmarcadas por la propuesta.
% fuente: 01_preguntas_picoc.md (origen de las preguntas y M8), 04_critica_abogado_diablo_cp1.md (ítem 8)

\subsection{Extracción, calidad y síntesis}
\label{sec:am-interna}
% PRISMA ítem 23c (limitaciones del proceso de revisión)

La doble extracción por IA alcanzó un acuerdo de \acuerdoNivelDos~\% en los campos codificados del Nivel 2, de modo que persiste error de extracción, que afecta sobre todo a los valores no reportados o no extraídos. El acuerdo por campo antes de la resolución fue bajo en la técnica de restauración del contexto (\acuerdoCampoTecnicaContexto{}), en los pilares centrales (\acuerdoCampoPilaresCentrales{}) y en la representación del estado (\acuerdoCampoRepresentacion{}), y mayor en la inferencia de \en{joins} (\acuerdoCampoInferenciaJoins{}), el mecanismo de PI4 (\acuerdoCampoMecanismo{}) y el historial de la evaluación (\acuerdoCampoHistorial{}). Los campos de bajo acuerdo son los de categorías más abiertas y los que definen los pilares centrales, de modo que las taxonomías de PI1 y PI2 y la clasificación por pilar son las más expuestas a error; el desempate y la verificación de la autora lo mitigan, pero no lo eliminan. Las verificaciones de la autora confirmaron todos los valores revisados (\verifC{}, \verifD{} y \verifB{}), lo que es compatible con una extracción correcta, pero también con sesgo de automatización o de confirmación. La muestra extraída a ciegas, de \nMuestraCiega{} estudios, es demasiado pequeña para estimar el acuerdo con precisión, y la codificación del historial, de la que depende la estratificación de PI5, se verificó de forma independiente solo en 3 de 5 estudios. El Nivel 1 se extrajo de las páginas iniciales, por lo que modelos y detalles se subestiman y las frecuencias «no reportado o no extraído» de un turno no son comparables con las multiturno.
% fuente: revision_4f.json (A5_acuerdo_campo_nivel2)

La IA propuso la puntuación de calidad (acuerdo de \acuerdoCalidad~\%), lo que puede clasificar mal los estudios cercanos al umbral de validez. Se mitigó con la revisión de los ítems de validez próximos al umbral y con la sensibilidad que reincorpora los estudios bajo el umbral. Los estudios de un turno no se evaluaron en calidad.

Los contrastes dentro del estudio están sujetos a sesgo de notificación, que tiende a sobreestimar los efectos. Algunas clases reúnen pocos estudios, y las clases «otros componentes» son heterogéneas. La confianza por hallazgo sigue reglas propias, no validadas, inspiradas en GRADE-CERQual, y las comparaciones entre estudios son solo descriptivas. Estas amenazas se mitigaron separando las clases, reportando la confianza, repitiendo el análisis sin \en{preprints}, contando los nulos con un umbral de $|\Delta| <$ \umbralNulo{}~pp y añadiendo una lectura con la confianza rebajada un nivel. Estas dos últimas son análisis \en{post hoc}.

\subsection{Validez externa y de conclusión}
\label{sec:am-externa}
% PRISMA ítems 23b (limitaciones de la evidencia) y 23c (limitaciones del proceso)

Solo se incluyeron estudios en inglés y español, y no se consultaron bases de datos en chino ni hispanas. La literatura publicada en otros idiomas queda fuera, lo que limita la generalización. La búsqueda se realizó el \fechaBusqueda{} y no se repitió, de modo que los estudios posteriores no están incluidos y el año 2026 es parcial, con tendencia truncada. La contaminación de \en{benchmarks} y el desfase de las líneas base pueden inflar los resultados reportados, y no se pudieron valorar. El análisis bibliométrico es exploratorio, y algunos estudios no se resolvieron en OpenAlex, por lo que depende de su cobertura.
